%\clearpage

\section {Introduction}
\label{sec:intro}

In this technical note, we report the results of a 3-flavour $\nu/\bar{\nu}$ joint analysis including the \nueCConepi{} sample,
performed by the {\em VALOR} group,        
on the combined Run 1-7c dataset, using data taken in both 
neutrino mode (\TotalPotNuRunOneSevenC{} POT),
with the polarity set to forward horn current (FHC),
and antineutrino (\TotalPotNuBarRunOneSevenC{} POT) mode
with polarity set reverse horn current (RHC).
The Run 7a-c data set was collected in the following periods: Run 7a \( \nu \) mode (\RunSevenAPOTOne{} POT, \RunSevenAPOTTwo{} POT) (\RunSevenADatesOne{},   \RunSevenADatesTwo{}), Run 7b \( \bar{\nu} \) mode (\RunSevenBPOTOne{} POT, \RunSevenBPOTTwo{} POT) (\RunSevenBDatesOne{},   \RunSevenBDatesTwo{} ) Run 7c \( \nu \) mode (\RunSevenCPOTOne{} POT) (\RunSevenCDatesOne{})

This analysis builds on our other analyses presented in Refs. 
\cite{OurRun1DisappAnalysis},
\cite{OurRun12DisappAnalysis}, 
\cite{OurRun12DisappAnalysis3Flavour},
\cite{OurRun123DisappAnalysis3Flavour},
\cite{OurRun1234aDisappAnalysis3Flavour},
\cite{OurRun1234DisappAnalysis3Flavour},
\cite{OurRun123JointOscillationAnalysis},
\cite{OurRun1234JointOscillationAnalysis},
\cite{OurRun56AntiNumuDisappearanceAnalysis3Flavour},
\cite{OurRun56AntiNueBarAppearanceAnalysis3Flavour},
and
\cite{OurRun1617b17cJointAnalysis3Flavour}.
In this analysis a hybrid frequentist-Bayesian method \cite{CousinsHighland} is employed, which consists of marginalising instead of profiling the nuisance parameters, as was used in the analyses \cite{OurRun56AntiNueBarAppearanceAnalysis3Flavour} and \cite{OurRun1617b17cJointAnalysis3Flavour}.

The consistency of the data with the expectation is given by comparing the predicted reconstructed neutrino energy 
and outgoing lepton angle spectra of 1 ring CCQE-like events and \nueCConepi{}-like events at Super-K to the observed spectra. 
The method of calculation of these predicted spectra is presented in \cref{sec:spectrum_prediction}.
This analysis uses flux tuning based on NA61/SHINE 2009 data \cite{FluxTuning}\cite{NA61/SHINET2K}. 
It also uses an indirect extrapolation from ND280 to Super-K, which is implemented via a fit of flux and cross-section parameters to
ND280 measurements of the muon momentum-angle distributions of $\bar{\nu}_\mu$ and  $\nu_\mu$
%CC0$\pi$-like, CC1$\pi$-like and CC other events whose 
for different event topologies
with the neutrino interaction vertex either in FGD1 or FGD2 (BANFF fit) \cite{BANFF2016}.


% Ereco Vs Theta mu-like and e-like
As in the previous analysis \cite{OurRun1617b17cJointAnalysis3Flavour},
the $e$-like samples, which now includes the \nueCConepi{} sample, have analysis templates that are built as a function of both 
the reconstructed neutrino energy and lepton angle,
in order to improve the separation of the different reaction / flavour components,
in particular $\bar{\nu}_{\mu} \rightarrow \bar{\nu}_e$ and $\nu_{\mu} \rightarrow \nu_e$,
fundamental in measuring $\delta_{CP}$.
Furthermore the $\nu_e$ intrinsic beam and NC components can also be better separated from signal.
The addition of the variable $\theta$ increases the number of template bins and so the CPU time,
the number of reconstructed neutrino energy bins has not been reduced,
in order to retain the previous sensitivity to oscillations.
For $\mu$-like samples the templates are built as a function of the reconstructed neutrino energy only
as in the past analyses.

Nominal predictions of the reconstructed neutrino energy and lepton angle spectrum 
of CC \( 0\pi \) events and \nueCConepi{}-like events in Super-K are presented for Run 1-7c in \cref{sec:spectrum_prediction}.

The oscillation analysis is performed using neutrino oscillation probabilities computed in a 3-flavour
framework including constant-density matter effects. 
Details of our oscillation probability calculations, estimates of the numerical accuracy of our
calculations and comparisons with calculations used by the alternative T2K oscillation analyses 
are given in Appendices A and B of Ref. \cite{OurRun123DisappAnalysis3Flavour}.

As in our previous analyses, the oscillation fit is performed using the likelihood-ratio method \cite{StatAnalysisCowan}.
The mathematical details of this method are given in Sec.\ 3.2 in Ref.\ \cite{OurRun1DisappAnalysis}.
We consider a total of \NSyst{} parameters including energy scale, neutrino flux, neutrino cross-section,
efficiency, initial-state nuclear environment and final-state re-scattering systematics 
(nuisance parameters) plus the oscillation parameters.
%In this Run 1-7c CC1pi analysis, 
All the nuisance parameters are marginalised~\cite{CousinsHighland}.
The nuisance parameters and their effects 
on the extrapolated single-ring event spectra are described in \cref{sec:systematics}.
The analysis method is presented in detail in \cref{sec:oscfit}.


Fit validations are presented in \cref{sec:validation}. 
The following updates have been added: 
the coverage studies study for the $\delta_{CP}$ vs mass hierarchy fit is shown in \cref{sec:coverage_deltacp}.
The main goal of this analysis is to search for CP violation in the context of the PMNS matrix
and new confidence intervals, produced with the Feldman-Cousins method~\cite{FeldmanCousins},
are shown in \cref{sec:fc_significance}.
These sections as well as appendix~\ref{app:method_significance} have been added in this new version.  

Confidence intervals of the other oscillation parameters are set using the 
constant $\Delta \chi^2$ method.
Sensitivity studies are shown in \cref{sec:sensi}.
The results of the measurement using the Run 1-7c dataset are given in \cref{sec:results} of this note. 

%In Sec.~\ref{sec:nuebarapp},~\ref{sec:sensitivity_nuebarapp} and~\ref{sec:results_nuebar} the $\bar{\nu}_e$ appearance analysis is shown 

%while the $\bar{\nu}_{\mu}$ disappearance analysis is shown in Sec.~\ref{sec:numubardis} and~\ref{sec:oscfit_numubardis}.
Finally, a summary is given in \cref{sec:summary}.  

\clearpage

%%% Local Variables:
%%% mode: latex
%%% TeX-master: "jointfit_5sample2016"
%%% End:
